% The title page carries [CONFIRM] items rather than \todo{} markers because
% todonotes places a float and floats are lost inside a titlepage.
\todo[inline]{CONFIRM: title with Dr.\ Zhang. Option~2 of three was accepted
pending supervisor confirmation; the alternatives and their trade-offs are in
the divergence analysis.}
\todo[inline]{CONFIRM: word limit and submission format. This draft is
approximately \num{29000} body words. If a 12{,}000--15{,}000 limit is
confirmed, the intended reduction is structural rather than editorial: move
\cref{ch:master} (master arm hardware and firmware), \cref{ch:development}
(development history) and \cref{ch:instrument} (the instrument chapter) to
appendices, which removes roughly \num{5500} words from the body while keeping
them examinable, and compress \cref{ch:results} by moving the per-measurement
detail into the evidence index of \cref{app:evidence}.}
\todo[inline]{FIGURE: Imperial College logo. Overleaf converts EPS
automatically, so restoring the \texttt{includegraphics} line in
\repo{title/title.tex} should be sufficient there; locally it needs
\texttt{epstopdf title/logo.eps} and ghostscript.}

\begin{abstract}
The base of a supernumerary robotic limb is a person. It sways, it is soft and
irreplaceable, and it cannot be told to hold still. This thesis reports a
teleoperation system for two seven-degree-of-freedom arms carried on a worn
backpack frame and commanded by a second person, together with what measuring
that system revealed about what it can and cannot do.

Two design premises shaped the work. The first is that commanding fourteen
joints through a body-worn master is a usability problem rather than a wiring
one, so position is taken from three or four well-conditioned channels,
orientation from an inertial sensor, and the remainder is supplied by autonomy
or left uncommanded. The second is that a master routing fourteen analog
channels through rotating joints will lose channels as connectors fatigue, so
losing them has to be survivable. Channel health is therefore classified per
channel, frame validation is scoped to the channels the active mode consumes,
and each channel can be disabled at run time with its cost published rather
than inferred. When no well-conditioned channel remains, the system reports
position unavailable instead of extrapolating one.

Position is decomposed so that each component comes from whichever sensor
observes it most directly: elevation from the gravity vector, azimuth from a
single shoulder joint, and radius from a forward-kinematic magnitude that
tolerates a dead roll channel. Measured against \num{20430} recorded frames,
the reduction from four sensed joints to one costs \SI{63.8}{\milli\metre} of
commanded tip error. Six of fourteen channels met the health criterion, and
classifying them correctly reversed two verdicts that had stood for weeks,
both of which came from computing statistics across transmitted rows rather
than distinct sensor updates. Replaying the failing channels with the
operator's input held constant showed them injecting \SI{54.5}{\milli\metre}
of commanded motion nobody asked for, against a master reach of
\SI{272}{\milli\metre}.

Vertical and fore-and-aft command track correctly on both arms. Lateral does
not, and a rotation-invariant test settles why: directions intended to be
\SI{90}{\degree} apart measure \SI{137.7}{\degree} apart, so the commanded
directions are not an orthogonal triad and no axis mapping can make them one.
Fourteen of fourteen injected faults were handled, two of them only after
correcting which observable a safety mechanism was keyed on.

Characterising the workspace produced a negative result with a quantified
mechanical remedy. Under the as-built mount the two arms share no reachable
cell of a sixty-three-cell frontal grid, separated by a \SI{0.25}{\metre} dead
band, which blocks inter-arm handover for geometric reasons that no control or
planning change can address. Searching for a working position in front of the
wearer's own chest returns nothing at any table height or distance, and
deleting the wearer's arms from the model moves the boundary by no more than
\SI{75}{\milli\metre}, which identifies the torso rather than the limbs as the
constraint. A buildable mount modification raising the shared count to
twenty-two of sixty-three is quantified and deliberately not applied, because
applying it invalidates the home pose and every clearance figure derived from
it.

Three capabilities were added after the interface was characterised and are
reported here for the first time. Motion generation was moved from a
per-joint step limit to a jerk-limited, phase-synchronised generator, which
removed \SI{51.3}{\milli\metre} of unrequested Cartesian excursion and brought
peak commanded joint speed from \num{12.53} times the manufacturer's declared
limit to \num{1.00}. A perception stage was added that measures the working
surface and the objects on it rather than being told their coordinates,
recovering six objects with every centre within \SI{1}{\milli\metre} and every
dimension within \SI{3}{\milli\metre}, and the resulting map is the first
thing in this project that a motion planner has ever consulted. And a
positioning error budget was assembled from the individual contributions,
giving a worst case of \SI{33.4}{\milli\metre} against a
\SI{30}{\milli\metre} capture requirement, with the largest remaining term
being a camera mounting transform that has never been measured.

No result in this thesis was validated on a physical arm carrying a load, and
no participant data has been collected. What stands between the system and the
two-person study it was designed around is enumerated in priority order, and
the first item is two soldered joints.
\end{abstract}

\renewcommand{\abstractname}{Acknowledgements}
\begin{abstract}
I thank my supervisor, Dr.\ Dandan Zhang, for the direction of this project
and for consistently pushing the work toward what could be measured rather
than what could be claimed. Much of what is useful in this thesis exists
because a result was questioned before it was written down.

I thank my co-supervisor, Dr.\ Etienne Burdet, for the framing that changed
the shape of the project. Treating the wearer and the operator as two people
with different stakes, rather than as one role, is what turned a control
problem into a question worth asking, and it is the reason the workspace and
safety chapters exist in the form they do.

I am grateful to the Human Robotics Group for access to the Dr.\ Octopus
platform, and to everyone who kept two Kinova arms and a backpack frame
available and working.
\todo{CONFIRM: name the technical staff, laboratory members and fellow
students who gave practical help, and anyone who lent equipment or time on the
platform. Only the two supervisors and the group are named here, because those
are the only contributions on record; the rest must be added by hand before
submission.}

This project spent a great deal of its time discovering that the instrument
was wrong rather than the system. I am thankful to everyone who responded to a
surprising number with scepticism instead of enthusiasm, which is the only
reason several of them are not in this document as findings.

Finally, I thank my family and friends for their support over the course of
the year.
\todo{CONFIRM: personalise the closing paragraph.}
\end{abstract}
